% GENERATED FILE. Edit canonical lessons, then run npm run generate:books.
% canonical-source-hash: fnv1a64:4402547a4b42d75a
% canonical-lessons: SA-C131-krsnah, SA-C131-dhusarah, SA-C131-patalah, SA-C131-spastah, SA-C131-sthulah

\chapter{Black, Grey, Pink, Clear, Thick}
\label{ch:sa-black-grey-pink-clear-thick}
\hlchaptermodality{131}

\begin{chapteropening}
\textbf{By the end of this chapter:} \emph{I can say black, grey, pink, clear and thick.}

Complete the last lesson of chapter 131: I can say black, grey, pink, clear and thick.
\end{chapteropening}

\section[\texorpdfstring{\sk{कृष्णः} (kṛṣṇaḥ)}{कृष्णः (kṛṣṇaḥ)}]{\sk{कृष्णः} (kṛṣṇaḥ) — black, dark}

\label{lesson:SA-C131-krsnah}



\begin{quote}
Before the new one: say the Sanskrit for weak, then the Sanskrit for clear.
\end{quote}

\subsection*{You'll want to know: \sk{कृष्णः}}
\textbf{\sk{कृष्णः}} — \emph{kṛṣṇaḥ} — \textquotedblleft{}black, dark\textquotedblright{}.

\begin{grammarlens}[title={What you've built}]
One more word for this chapter.
\end{grammarlens}

\subsection*{Your turn}
\begin{itemize}
\raggedright
  \item \emph{Say it:} \emph{kṛṣṇaḥ}
  \item \emph{Say it:} \emph{kṛṣṇaḥ}, once more
  \item \emph{Say it:} say \emph{kṛṣṇaḥ}
  \item \emph{Recall:} say the Sanskrit for weak, then the Sanskrit for clear, then say \emph{kṛṣṇaḥ} again
\end{itemize}

\begin{tcolorbox}[breakable,colback=teal!4,colframe=teal!35!black,title={Before you move on}]
What does \sk{कृष्णः} mean? (\textquotedblleft{}Black, dark\textquotedblright{}.) Say it once more.
\end{tcolorbox}
\section[\texorpdfstring{\sk{धूसरः} (dhūsaraḥ)}{धूसरः (dhūsaraḥ)}]{\sk{धूसरः} (dhūsaraḥ) — grey}

\label{lesson:SA-C131-dhusarah}



\begin{quote}
Before the new one: say the Sanskrit for clear, then the Sanskrit for black.
\end{quote}

\subsection*{You'll want to know: \sk{धूसरः}}
\textbf{\sk{धूसरः}} — \emph{dhūsaraḥ} — \textquotedblleft{}grey\textquotedblright{}.

\begin{grammarlens}[title={What you've built}]
One more word for this chapter.
\end{grammarlens}

\subsection*{Your turn}
\begin{itemize}
\raggedright
  \item \emph{Say it:} \emph{dhūsaraḥ}
  \item \emph{Say it:} \emph{dhūsaraḥ}, once more
  \item \emph{Say it:} say \emph{dhūsaraḥ}
  \item \emph{Recall:} say the Sanskrit for clear, then the Sanskrit for black, then say \emph{dhūsaraḥ} again
\end{itemize}

\begin{tcolorbox}[breakable,colback=teal!4,colframe=teal!35!black,title={Before you move on}]
What does \sk{धूसरः} mean? (\textquotedblleft{}Grey\textquotedblright{}.) Say it once more.
\end{tcolorbox}
\section[\texorpdfstring{\sk{पाटलः} (pāṭalaḥ)}{पाटलः (pāṭalaḥ)}]{\sk{पाटलः} (pāṭalaḥ) — pink}

\label{lesson:SA-C131-patalah}



\begin{quote}
Before the new one: say the Sanskrit for black, then the Sanskrit for grey.
\end{quote}

\subsection*{You'll want to know: \sk{पाटलः}}
\textbf{\sk{पाटलः}} — \emph{pāṭalaḥ} — \textquotedblleft{}pink\textquotedblright{}.

\begin{grammarlens}[title={What you've built}]
One more word for this chapter.
\end{grammarlens}

\subsection*{Your turn}
\begin{itemize}
\raggedright
  \item \emph{Say it:} \emph{pāṭalaḥ}
  \item \emph{Say it:} \emph{pāṭalaḥ}, once more
  \item \emph{Say it:} say \emph{pāṭalaḥ}
  \item \emph{Recall:} say the Sanskrit for black, then the Sanskrit for grey, then say \emph{pāṭalaḥ} again
\end{itemize}

\begin{tcolorbox}[breakable,colback=teal!4,colframe=teal!35!black,title={Before you move on}]
What does \sk{पाटलः} mean? (\textquotedblleft{}Pink\textquotedblright{}.) Say it once more.
\end{tcolorbox}
\section[\texorpdfstring{\sk{स्पष्टः} (spaṣṭaḥ)}{स्पष्टः (spaṣṭaḥ)}]{\sk{स्पष्टः} (spaṣṭaḥ) — clear}

\label{lesson:SA-C131-spastah}



\begin{quote}
Before the new one: say the Sanskrit for grey, then the Sanskrit for pink.
\end{quote}

\subsection*{You'll want to know: \sk{स्पष्टः}}
\textbf{\sk{स्पष्टः}} — \emph{spaṣṭaḥ} — \textquotedblleft{}clear\textquotedblright{}.

\begin{grammarlens}[title={What you've built}]
One more word for this chapter.
\end{grammarlens}

\subsection*{Your turn}
\begin{itemize}
\raggedright
  \item \emph{Say it:} \emph{spaṣṭaḥ}
  \item \emph{Say it:} \emph{spaṣṭaḥ}, once more
  \item \emph{Say it:} say \emph{spaṣṭaḥ}
  \item \emph{Recall:} say the Sanskrit for grey, then the Sanskrit for pink, then say \emph{spaṣṭaḥ} again
\end{itemize}

\begin{tcolorbox}[breakable,colback=teal!4,colframe=teal!35!black,title={Before you move on}]
What does \sk{स्पष्टः} mean? (\textquotedblleft{}Clear\textquotedblright{}.) Say it once more.
\end{tcolorbox}
\section[\texorpdfstring{\sk{स्थूलः} (sthūlaḥ)}{स्थूलः (sthūlaḥ)}]{\sk{स्थूलः} (sthūlaḥ) — thick, fat}

\label{lesson:SA-C131-sthulah}



\begin{quote}
Before the new one: say the Sanskrit for pink, then the Sanskrit for clear.
\end{quote}

\subsection*{You'll want to know: \sk{स्थूलः}}
\textbf{\sk{स्थूलः}} — \emph{sthūlaḥ} — \textquotedblleft{}thick, fat\textquotedblright{}.

\begin{grammarlens}[title={What you've built}]
One more word for this chapter.
\end{grammarlens}

\subsection*{Your turn}
\begin{itemize}
\raggedright
  \item \emph{Say it:} \emph{sthūlaḥ}
  \item \emph{Say it:} \emph{sthūlaḥ}, once more
  \item \emph{Say it:} say \emph{sthūlaḥ}
  \item \emph{Recall:} say the Sanskrit for pink, then the Sanskrit for clear, then say \emph{sthūlaḥ} again
\end{itemize}

\begin{tcolorbox}[breakable,colback=teal!4,colframe=teal!35!black,title={Before you move on}]
What does \sk{स्थूलः} mean? (\textquotedblleft{}Thick, fat\textquotedblright{}.) Say it once more.
\end{tcolorbox}
